
\documentclass[a4paper]{article}
\usepackage[pdftex]{graphicx}
\usepackage[utf8]{inputenc}
\usepackage{enumerate}
\usepackage{siunitx}
\sisetup{locale=DE}
\usepackage{href-ul}
\hypersetup{
	colorlinks=true,
	linkcolor=blue,
	urlcolor=blue} 
\usepackage{icomma}
\usepackage{amssymb}
\usepackage{geometry}
\geometry{a4paper, top=15mm, left=15mm, right=15mm, bottom=15mm,
	headsep=10mm, footskip=12mm}

\begin{document}
%Titel
{\bf Aufgaben zur Festigkeitslehre }\\
\begin{enumerate}[1.]

%Aufgabe 1
\begin{minipage}{0.48\textwidth}
\item Zur Untersuchung der Elastizität werden Probestäbe mit dem Querschnitt $A=\SI{50}{\milli\meter^2}$ durch Anhängen verschiedener Massen belastet und die dadurch verursachte Längenänderung mit einem Dehnungsmessstreifen gemessen. Für Stahl erhält man dabei nebenstehendes Diagramm.
\end{minipage}
\hfill
\begin{minipage}{0.48\textwidth}
\includegraphics[width=5 cm]{emodul219}
\end{minipage}

\begin{enumerate}[a)]
\item Bestimmen Sie mit Hilfe des gegebenen Diagramms
\begin{enumerate}[(i)]
\item den Elastizitätsmodul von Stahl
\item die Masse, bei welcher die elastische Formänderung $\varepsilon = \SI{10}{\micro\meter\over \meter}$ ist.
\item die Längenänderung in mm, die ein 2,5~m langer Stahlstab gleichen Querschnitts bei einer Spannung von $\SI{4}{\newton\milli\meter^{-2}}$ erfährt.
\end{enumerate}
\item Zeichnen Sie in das gegebene Diagramm den Graphen für einen Werkstoff mit einem halb so großen Elastizitätsmodul 
\end{enumerate}

%Aufgabe 2
\begin{minipage}{0.48\textwidth}
\item Ein Stahlrohr wird bei $\SI{20}{\celsius}$ spannungslos so eingebaut, dass es nicht nachgeben kann. Durch hindurchströmende heiße Gase wird es auf $\SI{80}{\celsius}$ erwärmt. 
\end{minipage}
\hspace{0.7 cm}
\begin{minipage}{0.48\textwidth}
\includegraphics[width=5 cm]{rohr219}
\end{minipage}

Der Längenausdehnungskoeffizient beträgt $ \alpha=\SI{12e-6}{\kelvin^{-1}}$, der Elastizitätsmodul $E=\SI{2,1e5}{\newton\milli\meter^{-2}}$
\begin{enumerate}[a)]
\item Berechnen Sie die dabei auftretende Druckspannung im Rohr.
\item Berechnen Sie die Druckkraft im Rohr
\end{enumerate} 

%Aufgabe 3
\begin{minipage}{0.55\textwidth}
\item Ein geteiltes Schwungrad wird durch zwei Schrumpfringe, die beide Nabenhälften umschließen, zusammengehalten. Die Streckgrenze des Materials beträgt $\SI{285}{\newton\milli\meter^{-2}}$, die Sicherheitszahl soll $2,0$ sein und die beiden Nabenhälften sollen mit der Kraft $F=\SI{1100}{\kilo\newton}$ zusammengehalten werden.
\end{minipage}
\hfill
\begin{minipage}{0.42\textwidth}
\includegraphics[width=5 cm]{nabe219}
\end{minipage}
\begin{enumerate}[a)]
\item Berechnen Sie die notwendige Breite $b$ der Ringe
\item Mit welchem Innendurchmesser müssen die Ringe hergestellt werden?
\end{enumerate}

%Aufgabe 4
\begin{minipage}{0.6\textwidth}
\item Eine Kupferbuchse und zwei starre Scheiben sind durch eine zentrale Stahlschraube 
miteinander verbunden. Bei der Temperatur $t_0=\SI{15}{\celsius}$ wird das System spielfrei, aber ohne Vorspannung montiert. Berechnen Sie die in Büchse und Schraube auftretenden Spannungen bei einer Erwärmung auf $\SI{215}{\celsius}$. \\
$[E_{Cu}=\SI{1,1e5}{\newton\milli\meter^{-2}}; \quad  E_{St}=\SI{2,1e5}{\newton\milli\meter^{-2}}; \quad 
A_{St}=\SI{58}{\milli\meter^2}; \quad \alpha_{Cu}=\SI{1,65e-5}{\kelvin^{-1}}; \quad d_{ACu}=\SI{26}{\milli\meter}; 
\quad d_{ICu}=\SI{20}{\milli\meter}]$
\end{minipage}
\hfill
\begin{minipage}{0.38\textwidth}
\includegraphics[width=4 cm]{stacu219}
\end{minipage}

%Aufgabe 5
\item Aus einem Stahlblech mit der Zugfestigkeit $\SI{340}{\newton\milli\meter^{-2}}<R_m<\SI{440}{\newton\milli\meter^{-2}}$ und einer Dicke von 6,0~mm sollen Scheiben mit einem Durchmesser von 36~mm ausgeschnitten werden. Die Scherspannung beträgt höchstens $80\%$ der Zugfestigkeit. Berechnen Sie die maximale Druckspannung, die im Stempel des Schneidwerkzeugs wirkt.

%Aufgabe 6 

\end{enumerate}
\fbox{
	\begin{minipage}{0.5\textwidth}
		Weitere Arbeitsblätter gibt es unter 
		
		\href{https://www.okuyakl.de}{www.okuyakl.de}
	\end{minipage}
	\hfill
	\begin{minipage}{0.4\textwidth}
		\includegraphics[width=1.5 cm]{../../viecher/zwe03}
		\includegraphics[width=2 cm]{../../viecher/afanticon1}
		
\end{minipage}}
%Ende
\end{document}